\documentclass[11pt,a4paper]{article}

\usepackage[utf8]{inputenc}
\usepackage[T1]{fontenc}
\usepackage[spanish]{babel}
\usepackage{geometry}
\usepackage{graphicx}
\usepackage{array}
\usepackage{tabularx}
\usepackage{makecell}
\usepackage[table]{xcolor}
\usepackage{fancyhdr}
\usepackage{titlesec}
\usepackage{helvet}
\usepackage{ragged2e}
\usepackage{enumitem}
\usepackage{setspace}
\usepackage{hyperref}
\usepackage{booktabs} % Para tablas con acabado profesional
\usepackage{enumitem} % Para control fino de las listas

\renewcommand{\familydefault}{\sfdefault}
\geometry{
  a4paper,
  left=3.0cm,
  right=3.0cm,
  top=2.35cm,
  bottom=2.35cm,
  headheight=1.9cm,
  headsep=0.45cm
}

\setlength{\parindent}{0pt}
\setlength{\parskip}{0.22em}
\setlength{\tabcolsep}{6pt}
\renewcommand{\arraystretch}{1.08}

\pagestyle{fancy}
\fancyhf{}
\fancyhead[L]{\includegraphics[height=1.35cm]{reporte_tecnico_overleaf_logo.png}}
\fancyhead[R]{\vspace{0.1cm}\small FEIRNNR - Carrera de Computación}
\renewcommand{\headrulewidth}{0.5pt}
\renewcommand{\footrulewidth}{0pt}
\makeatletter
\renewcommand{\headrule}{%
  {\color{gray!65}\hrule width\headwidth height\headrulewidth \vskip-\headrulewidth}%
}
\makeatother

\titleformat{\section}
  {\bfseries\fontsize{15}{17}\selectfont}
  {\thesection.}{0.8em}{}
\titlespacing*{\section}{0pt}{0.7ex}{0.2ex}

\newcolumntype{L}[1]{>{\RaggedRight\arraybackslash}p{#1}}
\newcolumntype{Y}{>{\RaggedRight\arraybackslash}X}

\begin{document}
\thispagestyle{fancy}

\vspace*{0.15cm}

\begin{center}
{\bfseries\fontsize{23}{26}\selectfont
Reporte Técnico de Actividades\\
Práctico-Experimentales Nro. 001\par}
\end{center}

\section{Datos de Identificación del Estudiante y la Práctica}

\noindent
\begin{tabularx}{\textwidth}{|L{0.43\textwidth}|Y|}
\hline
\textbf{Nombre del estudiante(s)} & Jaime Alejandro Landázuri Román \\ \hline
\textbf{Asignatura} & Simulación \\ \hline
\textbf{Ciclo} & 5 A \\ \hline
\textbf{Unidad} & 1 \\ \hline
\textbf{Resultado de aprendizaje de la unidad} & R1. Aplica análisis estadísticos de los datos
de la Simulación para mejorar la calidad de los resultados, bajo los principios de solidaridad, transparencia, responsabilidad y honestidad. \\ \hline
\textbf{Práctica Nro.} & 001 \\ \hline
\textbf{Título de la Práctica} & Construcción y simulación computacional de un modelo matemático \\ \hline
\textbf{Nombre del Docente} & José O. Guamán Q. \\ \hline
\textbf{Fecha} & 01 de Octubre de 2026 \\
 & 02 de Octubre 2026 \\ \hline
\textbf{Horario} & Jueves 10h30 a 11h30 \\
 & Viernes 07h30 a 09h30 \\ \hline
\textbf{Lugar} & Aula \\ \hline
\textbf{Tiempo planificado en el Sílabo} & 3 horas \\ \hline
\end{tabularx}

\section{Objetivo(s) de la Práctica}
Comprender el proceso de construcción de un modelo matemático y analizar su comportamiento mediante Python.
\begin{enumerate}
    \item Identificar variables, parámetros y relaciones matemáticas de un sistema.
    \item Implementar un modelo matemático en Python.
    \item Analizar gráficamente la influencia de los parámetros.
\end{enumerate}

\section{Materiales y reactivos}
\begin{itemize}
    \item NumPy.
    \item Matplotlib.
\end{itemize}

\section{Equipos y herramientas}
\begin{itemize}
    \item Computador.
    \item Python 3.
    \item Java.
    \item Editor de código.
\end{itemize}

\section{Procedimiento / Metodología Ejecutada}
Un modelo matemático representa un fenómeno mediante variables y relaciones matemáticas. Queremos simular el comportamiento de la atmósfera durante 24 horas y determinar en qué horas nuestro modelo considera que existe posibilidad de lluvia.
\centerline{$I=0.5H+0.3N+0.2T_f$}
Donde:
\begin{itemize}
    \item H = Humedad Normalizada.
    \item B = Nubosidad Normalizada.
    \item Tf = Factor de Temperatura.
\end{itemize}
Se requiere predecir en qué momento lloverá, para el factor de temperatura se tiene lo siguiente: \\
\centerline{Factor de Temperatura}
\begin{table}[h!]
\centering
\definecolor{headerblue}{RGB}{19, 93, 126}
\definecolor{rowlight}{RGB}{205, 235, 248}
\definecolor{borderblue}{RGB}{64, 160, 206}
\setlength{\arrayrulewidth}{0.75pt}
\arrayrulecolor{borderblue}
\renewcommand{\arraystretch}{1.25}
\begin{tabular}{|l|c|}
\hline
\rowcolor{headerblue}
\textbf{\textcolor{white}{Temperatura}} & \textbf{\textcolor{white}{T\textsubscript{f}}} \\ \hline
$\mathbf{\le 10\ ^\circ C}$ & 1.00 \\ \hline
\rowcolor{rowlight}
$\mathbf{12\ ^\circ C}$ & 0.90 \\ \hline
$\mathbf{14\ ^\circ C}$ & 0.80 \\ \hline
\rowcolor{rowlight}
$\mathbf{16\ ^\circ C}$ & 0.70 \\ \hline
$\mathbf{18\ ^\circ C}$ & 0.60 \\ \hline
\rowcolor{rowlight}
$\mathbf{20\ ^\circ C}$ & 0.50 \\ \hline
$\mathbf{22\ ^\circ C}$ & 0.40 \\ \hline
\rowcolor{rowlight}
$\mathbf{24\ ^\circ C}$ & 0.30 \\ \hline
$\mathbf{26\ ^\circ C}$ & 0.20 \\ \hline
\rowcolor{rowlight}
$\mathbf{\ge 28\ ^\circ C}$ & 0.10 \\ \hline
\end{tabular}
\end{table}

\centerline{Tabla de reglas}
\begin{table}[h!]
\centering
\definecolor{headerblue}{RGB}{19, 93, 126}
\definecolor{rowlight}{RGB}{205, 235, 248}
\definecolor{borderblue}{RGB}{64, 160, 206}
\setlength{\arrayrulewidth}{0.75pt}
\arrayrulecolor{borderblue}
\renewcommand{\arraystretch}{1.25}
\begin{tabular}{|l|l|}
\hline
\rowcolor{headerblue}
\textbf{\textcolor{white}{Índice I}} & \textbf{\textcolor{white}{Estado}} \\ \hline
\rowcolor{rowlight}
$\mathbf{I < 0.40}$ & Sin lluvia \\ \hline
$\mathbf{0.40 \le I < 0.60}$ & Baja posibilidad \\ \hline
\rowcolor{rowlight}
$\mathbf{0.60 \le I < 0.75}$ & Lluvia probable \\ \hline
$\mathbf{I \ge 0.75}$ & Lluvia \\ \hline
\end{tabular}
\end{table}

Ejemplo: \\
Supongamos que tenemos: \\
\begin{itemize}
    \item Humedad = 90%
    \item Nubosidad = 80%
    \item Temperatura = 18°C
\end{itemize}
Entonces \\
H = 0.9 \\
N = 0.8 \\
Tf = 0.6 \\
\\
Aplicamos $I=0.5H+0.3N+0.2T_f$ \\
\\
Calculamos $I=0.5(0.90)+0.3(0.80)+0.2(0.60)$ \\
$I = 0,81 = 81\%$ \\
Se debe programar en java y Python y llenar esta tabla además agregar las
graficas que se debe presentar \\
\begin{table}[h!]
\centering
\definecolor{headerblue}{RGB}{19, 93, 126}
\definecolor{rowlight}{RGB}{205, 235, 248}
\definecolor{borderblue}{RGB}{64, 160, 206}
\setlength{\arrayrulewidth}{0.75pt}
\arrayrulecolor{borderblue}
\renewcommand{\arraystretch}{1.25}
\begin{tabular}{|c|c|c|c|c|c|c|c|}
\hline
\rowcolor{headerblue}
\textbf{\textcolor{white}{Hora}} & \textbf{\textcolor{white}{Humedad}} & \textbf{\textcolor{white}{Nubosidad}} & \textbf{\textcolor{white}{Temp.}} & \textbf{\textcolor{white}{H}} & \textbf{\textcolor{white}{N}} & \textbf{\textcolor{white}{T\textsubscript{f}}} & \textbf{\textcolor{white}{Índice}} \\ \hline
\rowcolor{rowlight}
\textbf{06:00} & 65 & 40 & 14 & & & & \\ \hline
\textbf{08:00} & 70 & 50 & 16 & & & & \\ \hline
\rowcolor{rowlight}
\textbf{10:00} & 68 & 45 & 18 & & & & \\ \hline
\textbf{12:00} & 60 & 30 & 22 & & & & \\ \hline
\rowcolor{rowlight}
\textbf{14:00} & 75 & 70 & 20 & & & & \\ \hline
\textbf{16:00} & 85 & 85 & 18 & & & & \\ \hline
\rowcolor{rowlight}
\textbf{18:00} & 92 & 95 & 16 & & & & \\ \hline
\textbf{20:00} & 88 & 90 & 17 & & & & \\ \hline
\rowcolor{rowlight}
\textbf{22:00} & 80 & 75 & 15 & & & & \\ \hline
\end{tabular}
\end{table}

Ajuste el modelo y llene de nuevo la misma tabla además agregar las graficas que
se debe presentar

\section{Resultados esperados}
\begin{itemize}
    \item Modelo implementado correctamente.
    \item Gráficas de crecimiento y decrecimiento.
    \item Identificación de la influencia de r.
    \item Interpretación matemática del comportamiento.
\end{itemize}

\section{Preguntas de control}
\begin{enumerate}
    \item ¿Qué es un modelo matemático?
    \item ¿Cuál es la diferencia entre variable y parámetro?
    \item ¿Qué representa P₀?
    \item ¿Qué ocurre cuando r < 0?
    \item ¿Qué limitaciones tiene el modelo exponencial?
\end{enumerate}

\end{document}